
As we want to integrate over parts of a sphere, the simplest approach is to transform everything into spherical coordinates and compute the integrals that way.

In spherical coordinates, we have:
\begin{gather*}
x = r\sin\theta\cos\phi, \qquad
y = r\sin\theta\sin\phi, \qquad
z = r\cos\theta.
\end{gather*}

With $r=1$, this gives us:
\begin{gather*}
\omega =
\begin{bmatrix}
\sin\theta\cos\phi \\
\sin\theta\sin\phi \\
\cos\theta
\end{bmatrix}
\end{gather*}
and
\begin{gather*}
d\omega = \sin\theta \, d\theta \, d\phi.
\end{gather*}

\textbf{i.}

For the positive hemisphere, $\theta \in [0,\frac{\pi}{2}]$ and $\phi \in [0,2\pi]$:
\begin{gather*}
\int_{S^2_+} \cos\theta \, d\omega
= \int_{0}^{2\pi} \int_{0}^{\frac{\pi}{2}}
\cos\theta\sin\theta \, d\theta \, d\phi.
\end{gather*}

Here we can use $u$-substitution with $u=\sin\theta$ and $du=\cos\theta\,d\theta$:
\begin{gather*}
\int_{0}^{\frac{\pi}{2}} \cos\theta\sin\theta \, d\theta
= \int_{0}^{1} u \, du
= \frac{1}{2}.
\end{gather*}

And for the outer integral:
\begin{gather*}
\int_{0}^{2\pi} \frac{1}{2} \, d\phi = \pi.
\end{gather*}

\textbf{ii.}

For the entire sphere, $\theta \in [0,\pi]$ and $\phi \in [0,2\pi]$:
\begin{align*}
\int_{S^2} \omega \, d\omega
&= \int_{0}^{2\pi} \int_{0}^{\pi}
\begin{bmatrix}
\sin\theta\cos\phi \\
\sin\theta\sin\phi \\
\cos\theta
\end{bmatrix}
\sin\theta \, d\theta \, d\phi \\
&= \int_{0}^{2\pi} \int_{0}^{\pi}
\begin{bmatrix}
\sin^2\theta\cos\phi \\
\sin^2\theta\sin\phi \\
\cos\theta\sin\theta
\end{bmatrix}
\, d\theta \, d\phi.
\end{align*}

Here we can integrate each component separately.
The first two components will be zero, since:
\begin{gather*}
\int_{0}^{2\pi} \cos\phi \, d\phi = 0, \\
\int_{0}^{2\pi} \sin\phi \, d\phi = 0.
\end{gather*}

For the $z$ component, we can use the same $u$-substitution as in part (i):
\begin{gather*}
\int_{0}^{\pi} \cos\theta\sin\theta \, d\theta
= \int_{0}^{0} u \, du = 0.
\end{gather*}

This gives us:
\begin{gather*}
\boxed{
\int_{S^2} \omega \, d\omega =
\begin{bmatrix}
0 \\ 0 \\ 0
\end{bmatrix}
}
\end{gather*}

\textbf{iii.}

For the positive hemisphere, the $x$ and $y$ components will also be zero, as the integrals over $\phi$ remain unchanged.

For the $z$ component, we can again use $u$-substitution with $u=\sin\theta$ and $du=\cos\theta\,d\theta$:
\begin{gather*}
\int_{0}^{\frac{\pi}{2}} \cos\theta\sin\theta \, d\theta
= \int_{0}^{1} u \, du
= \frac{1}{2}.
\end{gather*}

And for the outer integral:
\begin{gather*}
\int_{0}^{2\pi} \frac{1}{2} \, d\phi = \pi.
\end{gather*}

This gives us:
\begin{gather*}
\boxed{
\int_{S^2_+} \omega \, d\omega =
\begin{bmatrix}
0 \\ 0 \\ \pi
\end{bmatrix}
}
\end{gather*}


\textbf{iv.}

For the negative hemisphere, the $x$ and $y$ components will again be zero, as the integrals over $\phi$ remain unchanged.

For the $z$ component, we integrate over
$\theta \in [\frac{\pi}{2},\pi]$ and $\phi \in [0,2\pi]$:
\begin{gather*}
\int_{0}^{2\pi}\int_{\frac{\pi}{2}}^{\pi}
\cos\theta\sin\theta \, d\theta \, d\phi.
\end{gather*}

Using $u$-substitution with $u=\sin\theta$ and
$du=\cos\theta\,d\theta$:
\begin{gather*}
\int_{\frac{\pi}{2}}^{\pi}
\cos\theta\sin\theta \, d\theta
= \int_{1}^{0} u \, du
= \left[\frac{u^2}{2}\right]_{1}^{0}
= -\frac{1}{2}.
\end{gather*}

And for the outer integral:
\begin{gather*}
\int_{0}^{2\pi} -\frac{1}{2} \, d\phi = -\pi.
\end{gather*}

This gives us:
\begin{gather*}
\boxed{
\int_{S^2_-} \omega \, d\omega =
\begin{bmatrix}
0 \\ 0 \\ -\pi
\end{bmatrix}
}
\end{gather*}
